\documentclass[a4paper,11pt]{article}
\usepackage[utf8]{inputenc}
\usepackage[T1]{fontenc}
\usepackage[ngerman]{babel}
\usepackage{booktabs}
\usepackage{tabularx}
\usepackage{array}
\usepackage{geometry}
\usepackage{parskip}
\usepackage{amssymb}

\geometry{margin=2.5cm}

\begin{document}
	
	\begin{center}
		\Huge{\textbf{RISC-V Befehlsübersicht}}
	\end{center}
	\vspace{12pt}
	In den Übungen zur Vorlesung \textit{Systemnahe Informatik} wird nur ein Teil des RISC-V-Befehlssatzes aktiv verwendet.
	Dabei handelt es sich um den RV32I-Basisbefehlssatz und die Multiplikations- und Divisionserweiterung RV32M.
	Das bedeutet, dass keine 64-Bit- und Gleitkommaarithmetik verwendet wird.
	Dadurch bleibt die Zahl der Befehle sehr überschaubar.
	
	% -----------------------------------------------------------------------
	\section{Unbedingte Sprungbefehle}
	
	Bei unbedingten Sprüngen wird der Programmzähler \texttt{pc} auf einen neuen Wert gesetzt.
	Um Rücksprünge zu ermöglichen, kann der vorherige Wert des \texttt{pc}-Registers im Register \texttt{ra} gesichert werden.
	
	\begin{center}
		\begin{tabular}{>{\ttfamily}lp{11cm}}
			\toprule
			\normalfont\textbf{Befehl} & \textbf{Beschreibung} \\
			\midrule
			j Label   & Springt zur Adresse \textit{Label}. \\
			jr Rs     & Indirekter Sprung zu der Adresse, die in Register \textit{Rs} steht. \\
			jal Label & Springt zur Adresse \textit{Label}.
						Der alte \texttt{pc} wird im Register \texttt{ra} als Rücksprungadresse gesichert.
						Dient dem Unterprogrammaufruf. \\
			jalr Rs   & Indirekter Sprung zu der Adresse, die in Register \textit{Rs} steht.
						Der alte \texttt{pc} wird im Register \texttt{ra} als Rücksprungadresse gesichert. \\
			ret       & Springe zur Adresse in Register \texttt{ra}. Dient dem Rücksprung aus einem Unterprogramm. \\
			\bottomrule
		\end{tabular}
	\end{center}
	
	% -----------------------------------------------------------------------
	\section{Bedingte Verzweigungsbefehle}
	
	Bei bedingten Sprüngen werden zuerst zwei Register miteinander verglichen.
	Trifft die kodierte	Bedingung zu, so wird der Sprung ausgeführt, andernfalls nicht.
	Es gibt keine Sprünge für die Bedingungen „kleiner gleich" und „echt größer als".
	Diese realisiert man, indem man	\texttt{bge} und \texttt{blt} verwendet, aber \textit{Rs} und \textit{Rt} vertauscht.
	
	Beim vorzeichenbehafteten Vergleich werden die beiden Register im Zweierkomplement interpretiert, d.\,h. sie stellen Zahlen von $-2\,147\,483\,648$ bis $+2\,147\,483\,647$ dar -- andernfalls werden sie als rein positive Zahlen im Bereich von $0$ bis $4\,294\,967\,295$ interpretiert.
	Es gibt keine Vergleiche mit Immediate-Konstanten, d.\,h. Konstanten müssen immer zuerst in ein	Register geladen werden.
	
	\begin{center}
		\begin{tabular}{>{\ttfamily}lp{7.5cm}}
			\toprule
			\normalfont\textbf{Befehl} & \textbf{Beschreibung} \\
			\midrule
			beq  Rs, Rt, Label & Springe, wenn $Rs = Rt$ \\
			bne  Rs, Rt, Label & Springe, wenn $Rs \neq Rt$ \\
			blt  Rs, Rt, Label & Springe, wenn $Rs < Rt$ (vorzeichenbehaftet) \\
			bge  Rs, Rt, Label & Springe, wenn $Rs \geq Rt$ (vorzeichenbehaftet) \\
			bltu Rs, Rt, Label & Springe, wenn $Rs < Rt$ (nur positive Zahlen) \\
			bgeu Rs, Rt, Label & Springe, wenn $Rs \geq Rt$ (nur positive Zahlen) \\
			\bottomrule
		\end{tabular}
	\end{center}
	
	\newpage
	
	% -----------------------------------------------------------------------
	\section{Arithmetische Befehle mit Konstante}
	
	Die folgenden Ganzzahl-Operationen verknüpfen den Inhalt eines Registers \textit{Rs} mit einem Im\-mediate-Wert \textit{Imm} im Bereich von $-2048$ bis $+2047$.
	Bei den Bitschiebebefehlen darf der Im\-medi\-ate-Wert nur zwischen $0$ und $31$ liegen.
	Das Ergebnis wird in ein weiteres Register \textit{Rd} geschrieben.
	
	\begin{center}
		\begin{tabular}{>{\ttfamily}lp{11cm}}
			\toprule
			\normalfont\textbf{Befehl} & \textbf{Beschreibung} \\
			\midrule
			add Rd, Rs, Imm & Addition \\
			slt Rd, Rs, Imm & Vorzeichenbehafteter Vergleich von \textit{Rs} und \textit{Imm}.
								Wenn \textit{Rs} kleiner ist, wird $1$ in \textit{Rd} geschrieben, andernfalls $0$. \\
			sltu Rd, Rs, Imm & Vergleich der vorzeichenlosen Werte von \textit{Rs} und \textit{Imm}.
								Wenn \textit{Rs} kleiner ist, wird $1$ in \textit{Rd} geschrieben, andernfalls $0$. \\
			and Rd, Rs, Imm & Bitweise Und-Verknüpfung \\
			or  Rd, Rs, Imm & Bitweise Oder-Verknüpfung \\
			xor Rd, Rs, Imm & Bitweise Exklusiv-Oder-Verknüpfung \\
			sll Rd, Rs, Imm & Bitweises Verschieben nach links, auffüllen mit Nullen \\
			srl Rd, Rs, Imm & Bitweises Verschieben nach rechts, auffüllen mit Nullen \\
			sra Rd, Rs, Imm & Bitweises Verschieben nach rechts, das Vorzeichenbit wird entsprechend kopiert. \\
			\bottomrule
		\end{tabular}
	\end{center}
	
	% -----------------------------------------------------------------------
	\section{Arithmetische Befehle auf Registern}
	
	Alle arithmetischen Befehle mit Konstanten können auch ausschließlich auf Register angewendet werden, d.\,h. der Inhalt von Register \textit{Rs} wird mit dem Inhalt von Register \textit{Rt} verknüpft und in Register \textit{Rd} geschrieben.
	Bei den bitweisen Verschiebungen werden nur die unteren 5\,Bit von \textit{Rt} als Verschiebungsweite interpretiert, die restlichen Bits werden ignoriert.
	
	Im Vergleich zu den arithmetischen Befehlen mit Konstante gibt es einige zusätzliche Verknüpfungen.
	Dies sind die Multiplikations- und Divisionsbefehle und die Subtraktion.
	Letztere hat keine Entsprechung mit unmittelbarer Konstante, da stattdessen eine negative Zahl addiert werden kann.
	
	\begin{center}
		\begin{tabular}{>{\ttfamily}lp{11cm}}
			\toprule
			\normalfont\textbf{Befehl} & \textbf{Beschreibung} \\
			\midrule
			add  Rd, Rs, Rt & Addition \\
			slt  Rd, Rs, Rt & Vorzeichenbehafteter Vergleich von \textit{Rs} und \textit{Rt}.
								Wenn \textit{Rs} kleiner ist, wird $1$ in \textit{Rd} geschrieben, andernfalls $0$. \\
			sltu Rd, Rs, Rt & Vergleich der vorzeichenlosen Werte von \textit{Rs} und \textit{Rt}.
								Wenn \textit{Rs} kleiner ist, wird $1$ in \textit{Rd} geschrieben, andernfalls $0$. \\
			and  Rd, Rs, Rt & Bitweise Und-Verknüpfung \\
			or   Rd, Rs, Rt & Bitweise Oder-Verknüpfung \\
			xor  Rd, Rs, Rt & Bitweise Exklusiv-Oder-Verknüpfung \\
			sll  Rd, Rs, Rt & Bitweises Verschieben nach links, auffüllen mit Nullen \\
			srl  Rd, Rs, Rt & Bitweises Verschieben nach rechts, auffüllen mit Nullen \\
			sra  Rd, Rs, Rt & Bitweises Verschieben nach rechts, das Vorzeichenbit wird entsprechend kopiert. \\
			sub  Rd, Rs, Rt & Subtraktion \\
			mul  Rd, Rs, Rt & Multiplikation \\
			div  Rd, Rs, Rt & Division (vorzeichenbehaftet) \\
			divu Rd, Rs, Rt & Division (nur positive Zahlen) \\
			rem  Rd, Rs, Rt & Divisionsrest (vorzeichenbehaftet) \\
			remu Rd, Rs, Rt & Divisionsrest (nur positive Zahlen) \\
			\bottomrule
		\end{tabular}
	\end{center}
	
	% -----------------------------------------------------------------------
	\section{Speicherzugriffe}
	
	Es gibt nur eine Adressierungsart: Register-indirekt mit Verschiebung.
	Das heißt, zur Berechnung der Speicheradresse wird eine Konstante \textit{Imm} (im Bereich von $-2048$ bis $+2047$) zum Inhalt des Registers \textit{Rs} addiert.
	Indem die Konstante auf $0$ gesetzt wird oder das Register \texttt{zero} verwendet wird, können damit auch register-indirekte und direkte Adressierung umgesetzt werden.
	
	\begin{center}
		\begin{tabular}{>{\ttfamily}lp{11cm}}
			\toprule
			\normalfont\textbf{Befehl} & \textbf{Beschreibung} \\
			\midrule
			lb  Rd, Imm(Rs) & Lade einen 8-Bit-Wert von der Adresse $Imm+Rs$ und erweitere den Wert entsprechend des Vorzeichens auf 32\,Bit, die in Register \textit{Rd} geschrieben werden. \\
			lh  Rd, Imm(Rs) & Lade einen 16-Bit-Wert von der Adresse $Imm+Rs$ und erweitere den Wert entsprechend des Vorzeichens auf 32\,Bit, die in Register \textit{Rd} geschrieben werden. \\
			lw  Rd, Imm(Rs) & Lade einen 32-Bit-Wert von der Adresse $Imm+Rs$ und schreibe ihn in Register \textit{Rd}. \\
			lbu Rd, Imm(Rs) & Lade einen 8-Bit-Wert von der Adresse $Imm+Rs$ und schreibe ihn in Register \textit{Rd}, wobei die restlichen Bits auf $0$ gesetzt werden. \\
			lhu Rd, Imm(Rs) & Lade einen 16-Bit-Wert von der Adresse $Imm+Rs$ und schreibe ihn in Register \textit{Rd}, wobei die restlichen Bits auf $0$ gesetzt werden. \\
			sb  Rt, Imm(Rs) & Schreibe die untersten 8\,Bit von \textit{Rt} an die Adresse $Imm+Rs$. \\
			sh  Rt, Imm(Rs) & Schreibe die untersten 16\,Bit von \textit{Rt} an die Adresse $Imm+Rs$. \\
			sw  Rt, Imm(Rs) & Schreibe den 32\,Bit breiten Inhalt von \textit{Rt} an die Adresse $Imm+Rs$. \\
			\bottomrule
		\end{tabular}
	\end{center}
	
	% -----------------------------------------------------------------------
	\section{Spezielle Befehle}
	
	Außer den hier aufgeführten Befehlen gibt es noch weitere, die aber im Rahmen der Vorlesung und Übung \textit{Systemnahe Informatik} keine Rolle spielen.
	
	\begin{center}
		\begin{tabular}{>{\ttfamily}lp{11cm}}
			\toprule
			\normalfont\textbf{Befehl} & \textbf{Beschreibung} \\
			\midrule
			nop          & No Operation: führe einen Takt lang nichts aus. \\
			scall        & Rufe eine Betriebssystemfunktion auf. Siehe Anhang~B. \\
			lui   Rd, ImmU & Schiebt \textit{ImmU} um 12\,Bit nach links (Multiplikation mit 4096) und schreibt diesen Wert in \textit{Rd}. \\
			auipc Rd, ImmU & Schiebt \textit{ImmU} um 12\,Bit nach links (Multiplikation mit 4096), addiert den Programmzähler \texttt{pc} und schreibt das Ergebnis in \textit{Rd}. \\
			\bottomrule
		\end{tabular}
	\end{center}
	
	% -----------------------------------------------------------------------
	\newpage
	\appendix
	
	\section{Registernamen}
	
	Grundsätzlich sind alle Register außer Register~0 (\texttt{zero}) gleichberechtigt, aber durch die ABI (Application Binary Interface) und die Calling Convention bekommen sie eine empfohlene Verwendung, die sich auch im Namen widerspiegelt.
	
	\begin{center}
		\begin{tabular}{clp{7cm}}
			\toprule
			\textbf{Nummer} & \textbf{Name} & \textbf{Beschreibung} \\
			\midrule
			0     & \texttt{zero}       & immer 0 \\
			1     & \texttt{ra}         & Rücksprungadresse \\
			2     & \texttt{sp}         & Stack Pointer \\
			3     & \texttt{gp}         & Global Pointer \\
			4     & \texttt{tp}         & Thread Pointer \\
			5--7  & \texttt{t0--t2}     & für temporäre Werte \\
			8     & \texttt{fp}         & Frame Pointer \\
			9     & \texttt{s1}         & zu sicherndes Register \\
			10--11 & \texttt{a0--a1}   & Rückgabewerte bzw. Funktionsparameter \\
			12--17 & \texttt{a2--a7}   & Funktionsparameter \\
			18--27 & \texttt{s2--s11}  & Register, die gesichert werden müssen \\
			28--31 & \texttt{t3--t6}   & für temporäre Werte \\
			\bottomrule
		\end{tabular}
	\end{center}
	
	% -----------------------------------------------------------------------
	\section{System-Calls}
	
	Im riscVivid-Simulator werden einige wenige System-Calls unterstützt, um mit dem Betriebssystem bzw. dem Benutzer in sehr eingeschränkter Weise zu kommunizieren.
	Ihre Verwendung entspricht der Linux-Spezifikation für den RISC-V.
	
	System-Calls werden durch die Instruktion \texttt{scall} aufgerufen.
	Die Nummer der gewünschten Funktion wird im Register \texttt{a7} gespeichert und eventuell nötige weitere Parameter in den Registern \texttt{a0} bis \texttt{a6}.
	
	\begin{center}
		\begin{tabular}{clp{2.5cm}p{7cm}}
			\toprule
			\textbf{a7} & \textbf{Linux-Name} & \textbf{Parameter} & \textbf{Beschreibung} \\
			\midrule
			63 & \texttt{read} & \texttt{a0=0} \newline \texttt{a1=Adresse} \newline \texttt{a2=Länge} &
				Lese eine Zeile vom Benutzer ein (stdin).
				Sie wird ab der Adresse \texttt{a1} abgelegt und nach \texttt{a2} Bytes abgeschnitten.
				Die Anzahl der tatsächlich eingegebenen Zeichen wird in Register \texttt{a0} zurückgeliefert. \\[6pt]
			64 & \texttt{write} & \texttt{a0=1} \newline \texttt{a1=Adresse} \newline \texttt{a2=Länge} &
				Gebe eine Zeichenkette aus (stdout).
				Die \texttt{a2} Bytes ab Adresse \texttt{a1} werden als Zeichenkette interpretiert und ausgegeben. \\[6pt]
			93 & \texttt{exit} & \texttt{a0=Exitcode} &
				Beende das Programm mit Exitcode \texttt{a0}. \\
			\bottomrule
		\end{tabular}
	\end{center}
	
\end{document}